% =============================================================================
% CAPÍTULO 1 — SPRINT 0: Incepción, Arquitectura Base y Planificación
% =============================================================================
\chapter{Sprint 0 — Incepción, Arquitectura Base y Planificación}
\label{cap:sprint0}
\resumenCapitulo{Fase de incepción del producto: establece la línea base tecnológica, la
arquitectura monolítica de despliegue, el \textit{Product Backlog} priorizado y la estrategia
de documentación continua que rigen los cinco incrementos posteriores de EthosHub.}

El Sprint 0, denominado estructuralmente como la fase de \textit{Inception}, constituye el
periodo de cimentación técnica y conceptual de la plataforma \textbf{EthosHub}. Durante este
ciclo, el equipo de \textbf{Byte Busters S.R.L.}\ ejecutó la transición de la visión de
negocio hacia una arquitectura de software desplegable y escalable, estableciendo la línea
base de configuración que regirá los cinco incrementos posteriores documentados en los
capítulos~\ref{cap:sprint1} a~\ref{cap:sprint5}.\par

\section{Objetivos del Sprint y Criterios de Aceptación Iniciales}
\label{sec:s0_objetivos}

El propósito primordial de este ciclo inicial es transformar la visión conceptual de
\textbf{EthosHub} en una arquitectura tangible y en un \textit{Product Backlog} priorizado.
Bajo el marco normativo de \textbf{ISO/IEC/IEEE 26515:2024}, el sprint no produce
funcionalidad de cara al usuario, sino los \textit{artefactos habilitadores} (arquitectura,
esquema de datos fundacional, andamiaje de repositorios y estrategia de documentación) que
permiten que el desarrollo sea escalable desde el primer \textit{commit}.\par

\textbf{Etimología y concepto:} el nombre combina el \textit{Ethos} —carácter e integridad
del profesional— con el \textit{Hub} —eje central de actividad—. Esta dualidad dicta que la
plataforma debe ser tan robusta en su \textit{backend} como cuidada en su \textit{frontend}.

\subsection{Meta de Negocio y Perfiles Objetivo}
\label{ssec:s0_meta}
La meta de negocio es posicionar a \textbf{EthosHub} como una plataforma \textit{SaaS} para
el sector tecnológico: no un generador de sitios estáticos, sino una herramienta de gestión
de marca personal y de descubrimiento de talento. El producto se diseña para tres perfiles
críticos:

\begin{itemize}[noitemsep]
    \item \textbf{Profesionales de TI:} construyen un portafolio público enriquecido
    (proyectos, experiencia, formación, \textit{skills} y CV en PDF) con una jerarquía visual
    clara.
    \item \textbf{Reclutadores:} descubren, filtran y conectan con talento mediante un
    \textit{pipeline} tipo Kanban, métricas e \textit{insights} asistidos por IA.
    \item \textbf{Evaluadores académicos:} validan la rigurosidad técnica y funcional del
    producto real en el marco del Taller de Ingeniería de Software (TIS).
\end{itemize}

\subsection{Objetivos Estratégicos del Sprint 0}
\label{ssec:s0_objetivos_estrategicos}
\begin{enumerate}[noitemsep]
    \item \textbf{Línea base tecnológica:} garantizar la paridad de entornos entre todos los
    desarrolladores, eliminando conflictos de versiones de JDK, Node y base de datos.
    \item \textbf{Arquitectura de identidad:} definir el núcleo del sistema en torno a la
    gestión de la cuenta y el perfil como primer punto de contacto del usuario.
    \item \textbf{Definición de alcance (\textit{grooming}):} transformar la visión de
    negocio en un \textit{Backlog} técnico priorizado de Épicas e Historias de Usuario.
    \item \textbf{Decisión arquitectónica de despliegue:} adoptar el modelo \textbf{monolítico},
    en el que la SPA de React se empaqueta dentro del JAR de Spring Boot (un único artefacto
    sirve la API y la interfaz).
    \item \textbf{Estrategia de documentación continua:} instaurar el paradigma
    \textit{Documentation as Code} exigido por ISO 26515.
\end{enumerate}

\subsection{Ficha Técnica del Sprint}
\label{ssec:s0_ficha}
La tabla~\ref{tab:s0_ficha} consolida los parámetros operativos del incremento. El criterio
de éxito transversal es la \textbf{validación de paridad de entornos al 100\,\%} en el equipo.

\begin{table}[!ht]
    \centering
    \renewcommand{\arraystretch}{1.4}
    \begin{tabularx}{\textwidth}{|l|X|}
        \hline
        \rowcolor{colorPrimario} \textcolor{white}{\textbf{Atributo}} & \textcolor{white}{\textbf{Detalle técnico}} \\ \hline
        \textbf{Periodo} & 16 al 28 de marzo de 2026 (12 días naturales). \\ \hline
        \textbf{Tag de cierre} & \comando{release/v0.1.0}. \\ \hline
        \textbf{Stack Frontend} & React 18.3 + TypeScript 5.3 / Vite 5.1 / Node.js 18+. \\ \hline
        \textbf{Stack Backend} & Java 17 / Spring Boot 4.0.6 / Maven 3.9+. \\ \hline
        \textbf{Base de datos} & PostgreSQL 15.10 autoalojada en el servidor TIS (schema \comando{core}, administrada vía phpPgAdmin). \\ \hline
        \textbf{Acceso a datos} & jOOQ (\comando{DSLContext}) + Spring JDBC. \\ \hline
        \textbf{Criterio de éxito} & Paridad de entornos validada al 100\,\% en el equipo. \\ \hline
    \end{tabularx}
    \caption{Ficha técnica resumida del Sprint 0 — EthosHub.}
    \label{tab:s0_ficha}
\end{table}

\begin{BreakingChange}
Cualquier desviación en la versión del JDK (\comando{Java 17}) o del entorno de ejecución de
Node.js (\comando{18+}) será tratada como un bloqueo crítico para la compilación del proyecto,
dado el empaquetado monolítico FE+BE descrito en la Sección~\ref{sec:s0_backend}.
\end{BreakingChange}

\subsection{Criterios de Aceptación Iniciales}
\label{ssec:s0_dod}
Para mantener el estándar corporativo se definieron criterios transversales obligatorios,
aplicables como \textit{Definition of Done} (DoD) a todas las Historias de los sprints
posteriores:

\begin{ChecklistDoD}
    \item \textbf{Validación en backend:} todo dato de entrada se valida con Jakarta Bean
    Validation, previniendo inyecciones de código.
    \item \textbf{Fidelidad de UI:} cumplimiento estricto de los \textit{mockups}, respetando
    la tipografía y la paleta corporativa.
    \item \textbf{Autenticación segura:} todo \textit{endpoint} privado se protege mediante
    filtros de seguridad y validación de tokens JWT.
    \item \textbf{Trazabilidad documental:} cada Historia cerrada actualiza la documentación
    técnica asociada en el mismo \textit{pull request}.
\end{ChecklistDoD}

\subsection{Hoja de Ruta de los Incrementos}
\label{ssec:s0_roadmap}
La planificación general (Plan General) distribuye el alcance en cinco incrementos de dos
semanas, según la tabla~\ref{tab:s0_roadmap}, tras la incepción del Sprint 0.

\begin{table}[!ht]
    \centering
    \renewcommand{\arraystretch}{1.3}
    \fontsize{9}{10.8}\selectfont\sffamily
    \begin{tabularx}{\textwidth}{|c|l|X|}
        \hline
        \rowcolor{colorPrimario}
        \textcolor{white}{\textbf{Sprint}} & \textcolor{white}{\textbf{Periodo}} & \textcolor{white}{\textbf{Módulo entregado}} \\ \hline
        1 & 30 mar -- 11 abr & Cuentas de Usuario y Perfil Personal. \\ \hline
        \rowcolor{grisTecnico}
        2 & 13 -- 25 abr & Habilidades y Validaciones de la Comunidad. \\ \hline
        3 & 27 abr -- 9 may & Vitrina de Proyectos y Conexión Externa. \\ \hline
        \rowcolor{grisTecnico}
        4 & 11 -- 23 may & Visibilidad, Marca Personal y Privacidad. \\ \hline
        5 & 25 may -- 6 jun & Estadísticas, Adaptabilidad y Lanzamiento. \\ \hline
    \end{tabularx}
    \caption{Hoja de ruta de los incrementos (Plan General).}
    \label{tab:s0_roadmap}
\end{table}

\clearpage
